\documentclass{amsart}
% For dealing with references we use the comment environment
\usepackage{verbatim}
\newenvironment{reference}{\comment}{\endcomment}
%\newenvironment{reference}{}{}
\newenvironment{slogan}{\comment}{\endcomment}
\newenvironment{history}{\comment}{\endcomment}

% For commutative diagrams we use Xy-pic
\usepackage[all]{xy}

% We use 2cell for 2-commutative diagrams.
\xyoption{2cell}
\UseAllTwocells

% We use multicol for the list of chapters between chapters
\usepackage{multicol}

% This is generally recommended for better output
\usepackage{lmodern}
\usepackage[T1]{fontenc}

% For cross-file-references

% Package for hypertext links:
\usepackage{hyperref}

% For any local file, say "hello.tex" you want to link to please

% Theorem environments.
%
\theoremstyle{plain}
\newtheorem{theorem}[subsection]{Theorem}
\newtheorem{proposition}[subsection]{Proposition}
\newtheorem{lemma}[subsection]{Lemma}

\theoremstyle{definition}
\newtheorem{definition}[subsection]{Definition}
\newtheorem{example}[subsection]{Example}
\newtheorem{exercise}[subsection]{Exercise}
\newtheorem{situation}[subsection]{Situation}

\theoremstyle{remark}
\newtheorem{remark}[subsection]{Remark}
\newtheorem{remarks}[subsection]{Remarks}

\numberwithin{equation}{subsection}

% Macros
%
\def\lim{\mathop{\mathrm{lim}}\nolimits}
\def\colim{\mathop{\mathrm{colim}}\nolimits}
\def\Spec{\mathop{\mathrm{Spec}}}
\def\Hom{\mathop{\mathrm{Hom}}\nolimits}
\def\Ext{\mathop{\mathrm{Ext}}\nolimits}
\def\SheafHom{\mathop{\mathcal{H}\!\mathit{om}}\nolimits}
\def\SheafExt{\mathop{\mathcal{E}\!\mathit{xt}}\nolimits}
\def\Sch{\mathit{Sch}}
\def\Mor{\mathop{\mathrm{Mor}}\nolimits}
\def\Ob{\mathop{\mathrm{Ob}}\nolimits}
\def\Sh{\mathop{\mathit{Sh}}\nolimits}
\def\NL{\mathop{N\!L}\nolimits}
\def\CH{\mathop{\mathrm{CH}}\nolimits}
\def\proetale{{pro\text{-}\acute{e}tale}}
\def\etale{{\acute{e}tale}}
\def\QCoh{\mathit{QCoh}}
\def\Ker{\mathop{\mathrm{Ker}}}
\def\Im{\mathop{\mathrm{Im}}}
\def\Coker{\mathop{\mathrm{Coker}}}
\def\Coim{\mathop{\mathrm{Coim}}}

% Boxtimes
%
\DeclareMathSymbol{\boxtimes}{\mathbin}{AMSa}{"02}

%
% Macros for moduli stacks/spaces
%
\def\QCohstack{\mathcal{QC}\!\mathit{oh}}
\def\Cohstack{\mathcal{C}\!\mathit{oh}}
\def\Spacesstack{\mathcal{S}\!\mathit{paces}}
\def\Quotfunctor{\mathrm{Quot}}
\def\Hilbfunctor{\mathrm{Hilb}}
\def\Curvesstack{\mathcal{C}\!\mathit{urves}}
\def\Polarizedstack{\mathcal{P}\!\mathit{olarized}}
\def\Complexesstack{\mathcal{C}\!\mathit{omplexes}}
% \Pic is the operator that assigns to X its picard group, usage \Pic(X)
% \Picardstack_{X/B} denotes the Picard stack of X over B
% \Picardfunctor_{X/B} denotes the Picard functor of X over B
\def\Pic{\mathop{\mathrm{Pic}}\nolimits}
\def\Picardstack{\mathcal{P}\!\mathit{ic}}
\def\Picardfunctor{\mathrm{Pic}}
\def\Deformationcategory{\mathcal{D}\!\mathit{ef}}

\setlength{\emergencystretch}{3em}
\begin{document}
\title[Stacks Project, Tag 0AAT]{713 --- Acyclic inverse limits over ordinals under surjective matching maps}
\maketitle
\noindent Copyright (C) 2005--2025 Johan de Jong. Stacks Project authors. Source: \href{https://stacks.math.columbia.edu/tag/0AAT}{Tag 0AAT}.

\noindent Editorial difficulty note (not author text). Difficulty H2: The proof uses transfinite recursion to lift cycles coherently, correcting primitive lifts by boundaries at successors and handling limit ordinals using the induction hypothesis. The ordinal matching-surjectivity condition is stronger than ordinary countable Mittag-Leffler arguments..

\noindent Editorial provenance note (not author text). History: excerpt from revision \texttt{a0444\allowbreak 6e57e\allowbreak c1fbc\allowbreak 252a8\allowbreak 71afc\allowbreak ec775\allowbreak 2fb28\allowbreak 07b14}, homology.tex, lines 7861--7913. Reference presentation was adapted; original-author context and core support blocks were retained or added. The mathematical wording is unchanged. Licensed under GNU FDL 1.2 or later; no invariant sections or cover texts. See ../licenses/COPYING. Context blocks reproduce author definitions and supporting results; they are not additional counted problems.

\begin{lemma}
\label{lemma-ML-over-ordinals}
Let $\alpha$ be an ordinal. Let $K_\beta^\bullet$, $\beta < \alpha$
be an inverse system of complexes of abelian groups over $\alpha$. If
for all $\beta < \alpha$ the complex $K_\beta^\bullet$ is acyclic and
the map
$$
K^n_\beta \longrightarrow \lim_{\gamma < \beta} K^n_\gamma
$$
is surjective, then the complex
$\lim_{\beta < \alpha} K_\beta^\bullet$ is acyclic.
\end{lemma}

\begin{proof}
By transfinite induction we prove this holds for every ordinal
$\alpha$ and every system as in the lemma. In particular, whilst
proving the result for $\alpha$ we may assume the complexes
$\lim_{\gamma < \beta} K^n_\gamma$ are acyclic.

\medskip\noindent
Let $x \in \lim_{\beta < \alpha} K^0_\beta$ with $\text{d}(x) = 0$.
We will find a $y \in \lim_{\beta < \alpha} K^{-1}_\beta$
with $\text{d}(y) = x$.
Write $x = (x_\beta)$ where $x_\beta \in K_\beta^0$ is the
image of $x$ for $\beta < \alpha$. We will construct $y = (y_\beta)$
by transfinite recursion.

\medskip\noindent
For $\beta = 0$ let $y_0 \in K_0^{-1}$
be any element with $\text{d}(y_0) = x_0$.

\medskip\noindent
For $\beta = \gamma + 1$ a successor, we have to find an element $y_\beta$
which maps both to $y_\gamma$ by the transition map
$f : K^\bullet_\beta \to K^\bullet_\gamma$ and to $x_\beta$ under the
differential. As a first approximation we choose $y'_\beta$ with
$\text{d}(y'_\beta) = x_\beta$. Then the difference $y_\gamma - f(y'_\beta)$
is in the kernel of the differential, hence equal to $\text{d}(z_\gamma)$
for some $z_\gamma \in K^{-2}_\gamma$.
By assumption, the map $f^{-2} : K^{-2}_\beta \to K^{-2}_\gamma$
is surjective. Hence we write $z_\gamma = f(z_\beta)$
and change $y'_\beta$ into $y_\beta = y'_\beta + \text{d}(z_\beta)$
which works.

\medskip\noindent
If $\beta$ is a limit ordinal, then we have the element
$(y_\gamma)_{\gamma < \beta}$ in $\lim_{\gamma < \beta} K^{-1}_\gamma$
whose differential is the image of $x_\beta$. Thus we can argue in exactly
the same manner as above using the termwise surjective map of complexes
$f : K_\beta^\bullet \to \lim_{\gamma < \beta} K_\gamma^\bullet$
and the fact (see first paragraph of proof) that we may assume
$\lim_{\gamma < \beta} K_\gamma^\bullet$ is acyclic by induction.
\end{proof}
\end{document}
